%%%%%%%%%%%%%%%%%%%%%%%%%%%%%%%%%%%%%%%%%%%%%%%%%%%%%
%			GENERÁTOR v. 18.09 				%
%%%%%%%%%%%%%%%%%%%%%%%%%%%%%%%%%%%%%%%%%%%%%%%%%%%%%
% Tento soubor slouží k jednotlivé zkusné kompilaci 
% písniček či k kompilaci samostatných pdf souborů 
% pomocí bashového skriptu kompilace.sh ve stejné 
% složce. Z důvodu hromadné kompilace musí poslední 
% dva řádky tohoto souboru zůstat nezměněny.
%%%%%%%%%%%%%%%%%%%%%%%%%%%%%%%%%%%%%%%%%%%%%%%%%%%%%
%			Jak kompilovat jednotlivé písně?        %
%%%%%%%%%%%%%%%%%%%%%%%%%%%%%%%%%%%%%%%%%%%%%%%%%%%%%
%	1. Do míst PÍSNIČKA upravte následnující řádek, 
%	   aby vypadal takto: 
%	   \input{../songy/[JMÉNO SOUBORU KOMPILOVNÉ PÍSNIČKY].tex}
%	2. Soubor písničky musí být přeložitelný a musí 
%	   se nacházet ve složce ../songy.
%%%%%%%%%%%%%%%%%%%%%%%%%%%%%%%%%%%%%%%%%%%%%%%%%%%%%
%			Jak psát soubory songů?                 %
%%%%%%%%%%%%%%%%%%%%%%%%%%%%%%%%%%%%%%%%%%%%%%%%%%%%%
%	1. Více návodu je k tomuto napsáno v souboru 
%      ../songy/00Songtemplate. 
%%%%%%%%%%%%%%%%%%%%%%%%%%%%%%%%%%%%%%%%%%%%%%%%%%%%%
%			Jak kompilovat celý zpěvník?			%
%%%%%%%%%%%%%%%%%%%%%%%%%%%%%%%%%%%%%%%%%%%%%%%%%%%%%
%	1. Více návodu je k tomuto napsáno v souboru
%	   ../Cely_zpevnik/zpevnik.tex.
%%%%%%%%%%%%%%%%%%%%%%%%%%%%%%%%%%%%%%%%%%%%%%%%%%%%%

% Hlavička START
\documentclass[openany,12pt]{memoir}
\usepackage[utf8]{inputenc} 
\usepackage[czech]{babel}
\usepackage[T1]{fontenc}
\usepackage[top=1.5cm, bottom=2cm, left=2cm, right=2cm]{geometry}  % --> NASTAVENÍ OKRAJŮ
\usepackage{fancyhdr}
\usepackage{graphicx}
\usepackage{xwatermark}
\usepackage{xcolor}
\usepackage{changepage}
\usepackage{pdfpages}
\usepackage{lettrine}
\usepackage{indentfirst}  %Důležité pro formátování

%%%%%%%%%%%%%%%%%%%%%%%%%%%%%%%%%%%%%%
%  FONT                              %
%%%%%%%%%%%%%%%%%%%%%%%%%%%%%%%%%%%%%%
\usepackage{amssymb}
\usepackage{tgschola}



%%%%%% Package na zpěvník
\usepackage[full]{leadsheets}%http://mirrors.nic.cz/tex-archive/macros/latex/contrib/leadsheets/leadsheets_en.pdf   --> dokumentace	
\definesongtitletemplate{empty}{} 
\setchords{
format = \bfseries \sffamily,   %tučné akordy
minor = {mi},% 
input-notation = {german},%
output-notation = {german}%
}
\definesongtitletemplate{empty}{} 


\newlength{\drop}
% VODOZNAK
\newwatermark[pages=1-,color=red!50,angle=0,scale=2, xpos=0,ypos=0]{\includegraphics[width=5cm]{obr/pozadi2.jpg}} %--> dvojka na pozadí


%%%%%%%%%%%%%%%%%%%%%%%%%%%%%%%%%%%%%%%%%%%%%%%%%%
%		 Vlastní příkazy
\newcounter{Slokočet}   %Automatické číslování slok
\newcommand{\mezera}{
\phantom{.}

}   %Horizontální odsazení slok (poněkud blbě zadefinovaný, ale jinak se formát rozbije jako wtf prostě)
\newcommand{\stred}{5.2cm}   %%% Na zarovnání slok doprostřed, pozn. automatičtější zarovnávání na střed nejde
\newcommand{\carka}{,\:}
\newcommand{\m}[1]{\color{white}{#1}}  %Pro akordy
\newcommand{\ap}{'}	%Pro apostrof
\newcommand{\elipsa}{\kern\fontdimen3\font} %Příkaz pro lepší zacházení s výpustkami (=...); je to vpodstatě jen mezera mezi tečkama výpustky
\newcommand{\pindent}{17.62482 pt} %Správná velikost \parindentu u layoutu se dvěma minipageama
\newcommand{\predtitle}{\huge}
\newcommand{\mezisloupci}{\phantom{TT}} %Místo mezi dvěma sloupci na jedné stránce
\newcommand{\z}{\hspace*{\fill}\null}

%%% Možné velikosti písem 
\newcommand{\normalni}{\normalsize}
\newcommand{\velky}{\fontsize{14.4}{15}\selectfont}
\newcommand{\vetsi}{\fontsize{15}{16}\selectfont}
\newcommand{\nejvetsi}{\fontsize{16}{17}\selectfont}
\newcommand{\nejnejvetsi}{\fontsize{17}{19}\selectfont}

%%% Stará definice sloky spoléhající na indenty
%\newlength{\pismeno}
%\settowidth{\pismeno}{x} %Tohle není moc ideální velikost, ale funguje
%\newif\ifslokavelka
%\slokavelkafalse
%\newcommand{\sloka}{
%\ifnum \value{Slokočet}>8  %Pokud je sloka dvouciferná
%\mezera \noindent \addtocounter{Slokočet}{1} \hspace*{-\pismeno}\arabic{Slokočet}.
%\else %Pokud jen jednociferná
%\mezera \noindent \addtocounter{Slokočet}{1} \arabic{Slokočet}. 
%\fi
%} 	%sloka, která se automaticky čísluje

\newcommand{\distanc}{\:}  %Vzdálenost čísla sloky před slokou
\newlength{\delkaargumentu}
%%% Sloka s automatickým číslováním
\newcommand{\sloka}{%
\addtocounter{Slokočet}{1}% Zvýší se o 1 počet slok
\mezera%  Sloka se odsadí vertikálně
\settowidth{\delkaargumentu}{\arabic{Slokočet}.\distanc}% Zde se určí délka odsazení 
\hspace*{-\delkaargumentu}%
\arabic{Slokočet}.\distanc%
\ignorespaces% Aby nevznikaly zbytečné mezery
}


%%% Sloka s vlastním argumentem
\newcommand{\ssloka}[1]{%     
\settowidth{\delkaargumentu}{#1\distanc}
\mezera%
\hspace*{-\delkaargumentu}%
#1\distanc%
\ignorespaces%
}  

%%% Refrén
\newcommand{\refren}[1][0]{%  Nepovinný argument sděluje, kolikátý refrén toto je, bez argumentu se vytiskne pouze refrén
\ifnum #1>0 %Pokud nepovinný argument existuje
\mezera%
\settowidth{\delkaargumentu}{\textbf{R$_{\text{#1}}$:}\distanc}%
\hspace*{-\delkaargumentu}%
\textbf{R$_{\text{#1}}$:}\distanc%
\ignorespaces%
\else %Pokud nepovinný argument neexistuje
\mezera%
\settowidth{\delkaargumentu}{\textbf{R:}\distanc}%
\hspace*{-\delkaargumentu}%
\textbf{R:}\distanc%
\ignorespaces%
\fi
}

\newcommand{\predehra}{\ssloka{\textbf{Předehra:}}}

%%% Capo
\newcommand{\kapodastr}[1]{%
\textit{Capo {#1}}
}


\addto\captionsczech{\renewcommand{\contentsname}{Seznam písní}}

%%%%%%%%%%%%%%%%%%%%%%%%%%%%%%%%%%%%
%    FORMÁTOVÁNÍ                   %
%%%%%%%%%%%%%%%%%%%%%%%%%%%%%%%%%%%%

%%% Vlevo zarovnaný text s blokem zarovnaným na střed
\usepackage{varwidth}% http://ctan.org/pkg/varwidth
\newenvironment{centerjustified}{%
  \begin{center} % so the minipage is centered
  \begin{varwidth}[t]{\textwidth}	
  \raggedright % so the minipage's text is left justified
  \setlength{\parindent}{\pindent}
}{%
  \end{varwidth}
  \end{center}
}


%Hlavnička END
\begin{document}
\sffamily %Sans-serif pro lepší čitelnost

\pagestyle{empty}
\newgeometry{top=1.5cm, bottom = 0cm, left = 2cm, right = 2cm}
	
\newpage

%Pro kompilaci je důležitý, aby ten \input a \end{document} zůstali na posledních
%dvou řádcích!!!!

\velky
\newpage
% PÍSNIČKA
\begin{song}{title=\predtitle\centering Za deset 10 \\\large Vlasta Redl  \vspace*{-0.3cm}}  %% sem se napíše jméno songu a autor
\begin{centerjustified}
\normalni

\begin{varwidth}[t]{0.50\textwidth}\setlength{\parindent}{0.0cm}  %Varianta č. 2 --> Dva sloupce

\sloka
	Ve středu ^{A \z}zabalím kufry a pojedu ^{Hmi \z}domů
	
	Stopnu si auto, co ^{E \z}zajíždí na ^{\z A}přehradu. ^{E}

	^{A \z}Pomůžu vyložit dynamit v ^{\z Hmi \,}kamenolomu.

	Pozdravím muže co ^{E}jsou tu za ^{\z A}velezradu.

	\mezera

	Ale tak udělám přece jen kolo, když už jsem tady.
	
	Je teprv poledne, a tak jsem celkem v plusu.

	Co je tu naboso šťastných krásných a mladých.

	Ale nejdřív si zjistím odjezdy autobusů

\refren
	Za deset ^{Hmi \z}deset ^{E} už nejede ^{A}nic. ^{F\#mi}

	Za deset ^{Hmi \z}deset ^{E} už nejede ^{A}nic.

\sloka
	Hned za hrází v jídelně u Mokrých sundám si oděv.

	Pochválím brýle nové brigádnici.

	Slíbím jí, že za ní vylezu po hromosvodě.
	
	Tak ahoj babi, zakončím pohlednici.
	
	Na dece na pláži vybalím Kamile, kde teďka hrajem.
	
	Natře mi záda vpodstatě jen za pusu.
	
	Přikreslí k adrese plánek celého kraje.
	
	Ještě zkontrolujeme si odjezdy autobusů.

\refren

\sloka
	Na Rudém právu počkám, až pytláci dopečou pstruhy.

	Přidám se k chlapům do války u Stalingradu.

	Kerouac Turgeněv Baudelaire Jan Pavel II.
	
	Básnicí všech zemí tak ještě pivo a jdu.
	
	Lacina s Grycem už u Bušů jen pro mě zastřelí prase.

	Přespolní vandrák vaří se zelím Led.
	
	Hej půjč mu to pádlo ať dá dítě v čase.
	
	Jsem hvězda a do první desky to mám deset let

\refren

\end{varwidth}\mezisloupci\begin{varwidth}[t]{0.55\textwidth}\setlength{\parindent}{0.05cm}
\vspace*{0.0525cm}  % V případě varianty č.2 jde odsud text do pravé části

\sloka
	V noci mi Martina vyzradí vojenské tajemství
	
	hlavního stanu Armády spásy.
	
	Je vdaná a večer vrací se z Bystřice zpět.
	
	A ráno mi Gábina s Vlaďkou rozčešou vlasy
	
	a než vypijem snídani je devět čtyřicet pět.
	
\refren

\sloka
	Poslední pivo s Evženem u šachu v zahrádce

	rušné

	sleduju, aby mi neujel zas.

	Vyměním dva pěšce za věž -- to je snad slušné

	a z koncových světel si odečtu svůj přesný čas.

	Za deset deset už nejede nic

	^{A}Ve středu zabalím kufr a zavolám ^{Hmi \z}taxík. ^{E A}

\end{varwidth}
\end{centerjustified}
\setcounter{Slokočet}{0}
\end{song}

\end{document}
